\begin{table}[t]
\centering
\small
\begin{tabular}{ccc cc cc c}
\toprule
& & & \multicolumn{2}{c}{Fixed SDSOS} & \multicolumn{2}{c}{SOS-basis SDSOS} & SOS \\
\cmidrule(lr){4-5}\cmidrule(lr){6-7}
\(n\) & \(d\) & \(m\) & Succ. & Gap & Succ. & Gap & Succ. \\
\midrule
2 & 2 & 6 & 0.60 & 1.3e+04 & 1.00 & 1.9e-07 & 1.00 \\
3 & 2 & 10 & 0.00 & -- & 0.90 & 1.5e-07 & 1.00 \\
4 & 2 & 15 & 0.00 & -- & 0.60 & 2.1e-07 & 1.00 \\
2 & 3 & 10 & 0.50 & 1.3e+05 & 0.90 & 1.6e-07 & 1.00 \\
3 & 3 & 20 & 0.00 & -- & 0.50 & 2.4e-07 & 1.00 \\
\bottomrule
\end{tabular}
\caption{
Initialization-bottleneck diagnostic on the generic SOS ensemble. 
Fixed SDSOS solves the problem in the monomial basis. 
SOS-basis SDSOS first solves the SOS relaxation, factors the SOS Gram matrix, and then re-solves SDSOS in the induced basis. 
This is an oracle diagnostic rather than a practical algorithm, since the basis is obtained from the SDP solution. 
The comparison tests whether SDSOS failure is partly caused by a poor initial basis.
}
\label{tab:initialization-bottleneck}
\end{table}
